\section{Introduction}

Fraud detection and related risk tasks at scale are still dominated by gradient-boosted decision
trees (GBDTs) over hand-engineered features~\cite{ke2017lightgbm,chen2016xgboost}. A compelling
alternative, mirroring foundation models elsewhere~\cite{bommasani2021foundation}, is to pretrain a
transformer on raw event sequences with a self-supervised objective and adapt it to downstream tasks
with a lightweight head --- a \emph{financial foundation model}
(FFM)~\cite{ostroukhov2026pragma,yeh2025treasure,padhi2021tabular}. Each transaction is tokenised
into field--value tokens; a set transformer encodes the fields of an event, and a temporal
transformer encodes an account's history of such events; adaptation reads the final-position
embedding with a linear probe or a light fine-tune.

In effect these models use \emph{two attentions}: one \emph{within} an event, over its fields, and
one \emph{across} an account's own history. Our observation is that this design sees only one entity
in isolation and is therefore \textbf{structurally blind to relational signal} --- patterns that
live in a shared entity's activity \emph{across} accounts. A compromised merchant produces a burst
of fraud spread over many cards; a fraud-farm IP drives many devices; a product goes viral across
many shoppers. None of these are visible to a model that attends only within a single account's
sequence, no matter how it is trained.

We add a \emph{third}, \textbf{cross-sequence attention}: for a target event at a shared entity
(merchant, IP, item, \dots), the model attends over that entity's recent events from \emph{other}
sequences, encoded by the same event encoder. Crucially we add it at \emph{fine-tuning time}, on top
of a pretrained two-attention backbone, rather than changing the pretraining recipe. This keeps the
FFM intact and makes the module a drop-in adaptation choice.

Our contributions are:
\begin{itemize}
\item \textbf{A third attention for FFMs (\S\ref{sec:method}).} A cross-sequence encoder that lets a
target event attend over a shared entity's recent cross-account events, added at fine-tuning time
and trained end-to-end with a linear head.
\item \textbf{Consistent relational gains (\S\ref{sec:results}).} Across three datasets where a
relational signal exists, the third attention improves \prauc{} by \num{+0.04} to \num{+0.07} over
standard fine-tuning of the same backbone.
\item \textbf{An isolation control.} On a synthetic benchmark with \emph{no} relational signal, the
gain \emph{collapses to zero} --- the same added capacity yields nothing --- establishing that the
improvement is the cross-sequence signal, not extra parameters.
\item \textbf{Beating a strong tabular baseline, and generalisation beyond fraud.} On TalkingData
the fine-tuned FFM with the third attention exceeds a strong GBDT; and on a non-fraud e-commerce
conversion task (Retailrocket) the same module delivers the same relational lift, showing the idea
is not fraud-specific.
\end{itemize}

All code, the synthetic generator, dataset adapters, and per-experiment result files are released for
reproducibility.
